\documentclass[noauthor,noinstructornotes]{ximera}

\title{Week 10}
\author{Math 1193 Design Team}

\begin{document}
\begin{abstract}
    Activities for Week 10.
\end{abstract}
\maketitle

\begin{problem}
Which of the following expressions are equal to 60? There are multiple
answers.

\begin{enumerate}
    \item \(2^2 \cdot 3 \cdot 5\)
    \item \(\frac{(2\cdot 3\cdot 5)^2}{15}\)
    \item \((60^2)^{1/2}\)
    \item \(2\cdot 3\cdot 5\cdot 2\)
    \item \(\frac{4\cdot 5}{3^{-1}}\)
    \item \(\frac{2 + 2^2 \cdot 3 \cdot 5}{2}\)
\end{enumerate}

What's the most convoluted expression equal to 60 you can come up
with using exponent rules?
\end{problem}

\begin{problem}
\begin{enumerate}
    \item What is \(\left(\frac{1}{3}\right)^0\)?
    \item What is \(\left(\frac{1}{3}\right)^1\)?
    \item What is \(\left(\frac{1}{3}\right)^{-1}\)?
    \item What is \(\left(\frac{1}{3}\right)^2\)?
    \item Is there a power of \(\frac{1}{3}\) that gives you a negative value?
    \item Draw the graph of \(f(x) =\left(\frac{1}{3}\right)^x\). Make sure to 
    include at least three labeled points.
    \item Describe in words the shape of the graph.
    \item Based on your graph, what are the domain and range of \(f\)?
    \item When will \(f(x) = 0\)?
    \item What is happening as \(x \to \infty\)? That is, as you plug
    bigger and bigger \(x\)-values into \(f\), what happens to the \(y\)-values?
\end{enumerate}
\end{problem}

\begin{problem}
\begin{enumerate}
    \item Find the \(y\)-intercept of the function defined by \(f\left(x\right)=7^x\).
    \item Find the \(y\)-intercept of the function defined by \(f\left(x\right)=2\cdot 7^x\).
    \item Find the \(y\)-intercept of the function defined by \(f\left(x\right)=\frac{1}{2}\cdot 7^x\).
    \item For a positive \(b\), find the \(y\)-intercept of the function defined by \(f\left(x\right)=a\cdot b^x\).
\end{enumerate}
\end{problem}

\begin{problem}
Lyneth has a chocolate bar weighing 4 ounces. Every day, 
she eats half of the chocolate remaining from the previous day.

\begin{enumerate}
    \item After 1 day, how much of the chocolate is left? 
    Give your answer in ounces.
    \item After 2 days?
    \item After 5 days? Feel free to leave your answer unsimplified.
    \item Will Lyneth ever run out of chocolate to eat if she
    continues to follow this process? Why or why not?
    \item Write an expression that gives the amount of chocolate left
    (in ounces) after \(x\) days.
\end{enumerate}

\begin{instructorNotes}
Use manipulatives if desired. You could cut up sheets of paper. This may
give a jumping-off point to a discussion of infinity and Zeno's paradox.
\end{instructorNotes}
\end{problem}

\begin{problem}
Using the rules of exponents, answer the following questions.
\textbf{Please do not look ahead until you have finished the current part!}

\begin{enumerate}
    \item Solve the equation \(2^1 \cdot 2^{-1} = 2^x\).
    \item Simplify \(2^1\) and \(2^{-1}\). Then multiply your answers together.
    What is the result?
    \item To summarize, we've shown that \(2^1 \cdot 2^{-1} = 2^0\), but
    also that \(2^1 \cdot 2^{-1} = 1\). What can you say about the value
    of \(2^0\)?
    \item Was there anything special about the base of 2 in this
    problem? What is \(17^0\)?
    \item (If you have time and energy) The expression \(0^0\) has 
    different interpretations depending on who you ask. The argument
    we used above to show that \(2^0 = 1\) does not work for \(0^0\).
    Why not?
\end{enumerate}
\end{problem}
\end{document}
